\chapter{Heat-kernel \(a_4\) и нормировка gauge coupling}

Локальный product lift свёл gauge mass к
\begin{equation}
 \frac{m_A^2}{\chi^2}=8g^2.
\end{equation}
Остаётся проверить, фиксирует ли тот же heat-kernel coefficient, который
дал scalar-нормировку \(\kappa=2\), значение \(g\).

\section{Laplace-type decomposition}

Для covariant product Dirac-оператора
\begin{equation}
 \mathcal D_A
 =i\gamma^\mu(\partial_\mu+A_\mu)
 +\gamma^5\Phi
\end{equation}
используем antihermitian convention
\begin{equation}
 A_\mu=iQ\,a_\mu,
 \qquad
 F_{\mu\nu}=iQ\,f_{\mu\nu}.
\end{equation}
В плоском фоне relevant endomorphism имеет блоки
\begin{equation}
 E_F=-\frac12\gamma^{\mu\nu}F_{\mu\nu},
 \qquad
 E_\Phi=-i\gamma^\mu\gamma^5D_\mu\Phi-\Phi^2,
\end{equation}
а curvature covariant derivative равна
\begin{equation}
 \Omega_{\mu\nu}=F_{\mu\nu}.
\end{equation}
Локальная часть \(a_4\) содержит
\begin{equation}
 \frac1{(4\pi)^2}
 \Tr\left(
 \frac12E^2+\frac1{12}\Omega_{\mu\nu}\Omega^{\mu\nu}
 \right).
\end{equation}

\section{Clifford traces}

При
\begin{equation}
 \operatorname{tr}_{\mathrm{spin}}\mathbf1=4
\end{equation}
и стандартной евклидовой Clifford-конвенции получаем
\begin{align}
 \operatorname{tr}_{\mathrm{spin}}
 \left(\frac12E_\Phi^2\right)
 &\supset
 2\,\Tr_{\mathrm{orb}}(D_\mu\Phi D^\mu\Phi),
 \\
 \operatorname{tr}_{\mathrm{spin}}
 \left(
 \frac12E_F^2+\frac1{12}\Omega^2
 \right)
 &=
 -\frac23\Tr_{\mathrm{orb}}(F_{\mu\nu}F^{\mu\nu}).
\end{align}
Так как \(F=iQf\), второй член положителен:
\begin{equation}
 -\Tr_{\mathrm{orb}}F^2
 =\Tr_{\mathrm{orb}}Q^2\,f^2.
\end{equation}
После деления всего \(a_4\)-блока на общий scalar coefficient \(2\)
нормированная лагранжева плотность содержит
\begin{equation}
 \mathcal L_4
 \supset
 \Tr_{\mathrm{orb}}(D_\mu\Phi D^\mu\Phi)
 +\frac13\Tr_{\mathrm{orb}}Q^2\,
 f_{\mu\nu}f^{\mu\nu}.
\end{equation}

\section{Orbit charge trace}

На минимальном reality-квадрате
\begin{equation}
 Q=\operatorname{diag}(0,1,-1,0),
\end{equation}
поэтому
\begin{equation}
 \Tr_{\mathrm{orb}}Q^2=2.
\end{equation}
Gauge kinetic term принимает вид
\begin{equation}
 \frac23 f_{\mu\nu}f^{\mu\nu}.
\end{equation}
Сопоставление с конвенцией
\begin{equation}
 \frac1{4g^2}f_{\mu\nu}f^{\mu\nu}
\end{equation}
даёт
\begin{equation}
 \boxed{g^2=\frac38,}
 \qquad
 \boxed{g=\sqrt{\frac38}.}
\end{equation}

\begin{theorem}[Common-\(a_4\) coupling pass]
В минимальном relative-\(U(1)\) product triple gauge coupling не является
новым непрерывным параметром после выбора общей \(a_4\)-нормировки.
Scalar и gauge kinetic terms фиксируют
\(g^2=3/8\) на spectral matching scale.
\end{theorem}

\begin{remark}
Утверждение относится к четырёхмерному zero-mode effective block.
Его перенос на полную \(M_4\times K\)-компактификацию требует отдельного
размерностного гейта.
\end{remark}

\section{Исправленный gauge mass и supertrace}

Подстановка найденного coupling в локальный product lift даёт
\begin{equation}
 \frac{m_A^2}{\chi^2}
 =8g^2=3.
\end{equation}
Gauge, Goldstone и ghost block поэтому добавляет
\begin{equation}
 \Delta N_{\mathrm{gauge+ghost}}
 =3\left(\frac{m_A^2}{\chi^2}\right)^2
 =3\cdot3^2=27.
\end{equation}
Совместно с finite scalar--fermion numerator \(40\) получаем
\begin{equation}
 N_0^{\mathrm{fin+gauge}}=40+27=67
\end{equation}
и
\begin{equation}
 \boxed{
 B_0^{\mathrm{fin+gauge}}
 =\frac{67}{64\pi^2}>0.}
\end{equation}

\begin{proposition}[Zero-mode completion milestone]
После включения выведенного relative-\(U(1)\) BV-блока полный
недилатонный zero-mode coefficient остаётся строго положительным. Его
значение \(67/(64\pi^2)\) не содержит свободного gauge coupling.
\end{proposition}

\section{Нормировочный red-team}

Общий spectral moment \(f_0\) умножает scalar и gauge части одинаково и
сокращается при фиксации канонической scalar-нормировки. Полный trace на
KO6-удвоении также удваивает одновременно
\(\Tr(D\Phi)^2\) и \(\Tr Q^2\), поэтому отношение и значение \(g\) не
меняются. Результат изменится только при:
\begin{enumerate}
 \item ином charge representation;
 \item независимой gauge kinetic добавке вне spectral action;
 \item иной нормировке генератора с изменённым глобальным периодом.
\end{enumerate}
Все три изменения означают новую версию parent model.

\begin{nogocriterion}[Gauge normalization kill-rule]
Если полный fluctuated Dirac trace содержит дополнительные charged
представления, меняющие \(\Tr Q^2=2\), они должны быть объявлены до
KK-аудита. Сохранение \(g^2=3/8\) при скрытом расширении charge ledger
запрещено.
\end{nogocriterion}

\begin{protocol}[Следующий гейт]
Zero-mode gauge completion закрыт. Теперь требуется перечислить spin и
flat-character ветви fermion operator на
\(K=\mathbb{RP}^3\times S^1\), вычислить для каждой ветви regulated
\(\Delta N_{\mathrm{KK}}\) в одной heat-kernel/дзета-схеме и проверить
устойчивость знака
\begin{equation}
 B_{\mathrm{full}}
 =\frac{67+c_\sigma^2+\Delta N_{\mathrm{KK}}}{64\pi^2}.
\end{equation}
\end{protocol}

Машинный аудит сохранён в
\path{s2t_v3_a4_gauge_coupling_results.json}.